\chapter{Stratégies mixtes}
Un des problèmes inhérents au concept d'équilibre de Nash en
stratégies pures est que pour certains jeux, de tels équilibres
n'existent pas. P.ex.\,le jeu de ``Pierre, Papier, Ciseaux" :
\begin{center}
\begin{game}[Pierre, Papier, Ciseaux]{3}{3}
     &  $Pi$    &  $Pa$  & $Ci$\\
$Pi$ & $0,0$    & $-1,1$ & $1,-1$\\
$Pa$ & $1,-1$   & $0,0$  & $-1,1$\\
$Ci$ & $-1,1$   & $1,-1$ & $0,0$\\
\end{game}\hspace*{20mm}%
\end{center}
n'admet pas d'équilibre de Nash.

La raison pour laquelle on n'a pas d'équilibres est la suivante :
la notion d'équilibre de Nash en stratégies pures suppose que
chaque joueur connaît les stratégies des autres joueurs. Or, nous
sommes dans des jeux ou chaque joueur a intérêt à cacher sa
stratégie, ou a bluffer. En effet, dans les jeux ``pile ou face"
ou ``tirer un penalty", ou ``bluff au poker", on n'utilise pas
toujours la même stratégie, et on ne connaît pas non plus à
l'avance celle de l'adversaire.
\begin{center}
\begin{game}[Pile ou face]{2}{2}
     &  $Pi$  &  $Fa$\\
$Pi$ & $1,-1$ & $-1,1$\\
$Fa$ & $-1,1$ & $1,-1$
\end{game}\hspace*{20mm}%
\end{center}

Les stratégies mixtes, ou aléatoires, vont permettre de
représenter ces possibilités de bluff, ou de jouer aléatoirement.

\section{Définition des concepts}
Nous allons définir les stratégies mixtes, puis les jeux dans
lesquels les joueurs ont la possibilité de jouer des stratégies
mixtes, et enfin définir un équilibre de Nash en stratégies mixtes
comme un équilibre de Nash d'un tel jeu.

\subsection{Stratégies mixtes}
Nous commençons par définir le concept de stratégies mixtes. Soit
donc $G = (I, (S_i)_i, (g_i)_i)$ un jeu sous forme normale

\begin{definition} Une strat\'egie mixte pour le joueur
$i$ est une loi de probabilit\'es sur $S_i$. $\Sigma_i
=\Delta(S_i)$ repr\'esente l'ensemble des strat\'egies
mixtes du joueur $i$. %
\end{definition}

On utilisera les notations : $\Sigma = \Pi_i \Sigma_i$,
$\Sigma_{-i} = \Pi_{j\neq i} \Sigma_j$.

Par opposition aux stratégies mixtes, on appelle les éléments de
$S_i$ {\bf stratégies pures}. Ce sont des stratégies
déterministes.

On utilisera selon les cas différentes notations pour une
stratégie mixte.
\begin{enumerate}
\item Notation fonction. $\sigma_i$ est l'application de $S_i$
vers $\reals$ qui associe à la stratégie pure $s_i$ sa probabilité
d'être jouée. Exemple: $S_1=\{H,B\}$,
$\sigma_1(H)=\sigma_1(B)=\frac12$%
\item Vecteur. On écrit $\sigma_i$ comme un vecteur de
probabilités. Si $S_i= \{s_{i,1},\ldots,s_{i,n_i}\}$, on
$\sigma_i$ se représente comme
$(\sigma_i(s_{i,1}),\ldots,\sigma_i(s_{i,n_i}))$. Exemple :
$\sigma_1=(\frac12,\frac12)$. %
\item Combinaison convexe de stratégies pures : $\frac12H +
\frac12B$.
\end{enumerate}

Il y a plusieurs interprétations possibles des stratégies mixtes,
parmi elles, les plus courantes sont les suivantes :
\begin{itemize}
\item Possibilit\'e ``mat\'erielle'' de jouer al\'eatoirement.
Rien ne s'oppose à priori à ce qu'un joueur décide aléatoirement
quelle stratégie pure utiliser. La stratégie mixte est un choix
aléatoire d'une stratégie pure.%
\item Bluff, et part d'incertitude que chacun laisse sur sa
strat\'egie. Ici, une stratégie mixte représente plutôt une
croyance que chaque joueur a sur les manières de jouer des autres
joueurs. Même si je décide si je vais bluffer ou non, de manière
pas forcément aléatoire, je ne souhaite surtout pas que les
autres joueurs connaissent mon choix !%
\item Dans une population de joueurs, une strat\'egie
correspond à une proportion dans laquelle une stratégie pure est
jouée dans la population. Par exemple, si dans la ville de New
Delhi 30\% des piétons cherchent à se ranger à gauche
(systématiquement) lorsqu'ils croisent un autre piéton, et 70\% à
droite, à chaque fois que je croise un piéton dans cette ville je
joue face à la stratégie mixte $(30\%,70\%)$.%%
\item Incertitude sur les paiements. (Harsanyi). Supposons que
selon mon ``humeur'' du moment, j'aie une faible préférence pour
bluffer, ou pour ne pas bluffer, et que les autres joueurs ne
connaissent pas mes préférences. Dans ce cas, je joue la stratégie
que je préfère, donc pas aléatoirement, mais les autres joueurs
ont l'impression que je joue aléatoirement, et se représentent mon
comportement par une stratégie mixte.
\end{itemize}

\begin{proposition}
L'ensemble $\Sigma_i$ des stratégies mixtes du joueur $i$ est
convexe. Ses points extr\^emaux sont les stratégies qui mettent
probabilité 1 sur
un seul point de $S_i$. %
\end{proposition}

Remarquons qu'un strat\'egie pure $s_i$ correspond \`a la
strat\'egie mixte qui joue $s_i$ avec probabilit\'e~1. On
considère donc $S_i$ comme sous-ensemble de $\Sigma_i$. Par
conséquent, toute stratégie pure est vue comme une stratégie
mixte. En considérant l'ensemble des stratégies mixtes, on étend
les ensembles des stratégies des joueurs.

\subsection{Extension mixte d'un jeu}

Si chaque joueur $j$ joue la strat\'egie $\sigma_j$, la
probabilité que $(s_j)_j$ soit le profil d'actions effectivement
joué est $\Pi_i \sigma_i(s_i)$. Par conséquence, le paiement
esp\'er\'e du joueur $i$ est :
$$
\sum_{(s_j)\in \prod_j S_j} \Pi_j \sigma_j(s_j) g_i((s_j)_j)
$$
Remarque : On suppose que les fonctions de gain $g_i\colon
S\to\reals$ sont des fonctions d'utilité de von Neumann et
Morgenstern, l'utilité pour l'aléa est l'espérance de l'utilité
obtenue.

La relation précédente définit des fonctions de paiement
$g_i\colon\Sigma\to\reals$. On utilise donc la même notation pour
l'application de $S$ vers $\reals$ et pour son extension de
$\Sigma$ vers $\reals$.

\begin{definition}
{\em L'extension mixte} du jeu sous forme normale $(I,
(S_i),(g_i)_i)$ est le jeu sous forme normale
$(I,(\Sigma_i)_i,(g_i)_i)$.%
\end{definition}

Dans l'extension mixte :
\begin{itemize}%
\item L'ensemble des joueurs est $I$.%
\item Chaque joueur choisit une stratégie mixte $\sigma_i\in
\Sigma_i$. %
\item Le paiement de $i$ est $g_i(\sigma)$.%
\end{itemize}%


\begin{proposition} L'application $g\colon\Sigma\to\reals$ est multilin\'eaire,
c.a.d.\, elle est lin\'eaire sur chaque coordonn\'ee $\Sigma_i$.
\end{proposition}
{\bf Preuve :} il faut montrer que pour
$\sigma_{-j}\in\Sigma_{-j}$, $\sigma_j, \sigma'_j \in \Sigma_j$ et
$\lambda\in[0,1]$, et $\sigma''_j=\lambda \sigma_j + (1-\lambda)
\sigma'_j$:
$$
g(\sigma_{-j},\sigma''_{j}) = \lambda g(\sigma_{-j},\sigma_{j}) +
(1-\lambda)g(\sigma_{-j},\sigma'_{j})
$$
En particulier on a :
$$
g(\sigma_{-j},\sigma_{j}) = \sum_{s_j\in S_j}
\sigma_j(s_j)g(\sigma_{-j},s_j)
$$

\subsection{Équilibre de Nash en stratégies mixtes}

\begin{definition}
 Un \'equilibre de Nash en strat\'egies mixtes
de $G$ est un \'equilibre de Nash du jeu avec ensembles de
strat\'egies $\Sigma_i$ et fonction de paiements $g_i$.
\end{definition}

Un équilibre de Nash en stratégies mixtes est donc un profil de
strat\'egies mixtes $\sigma\in\Sigma$ tel que $\forall i$,
$\forall \sigma'_i\in \Sigma_i$,
$$
g_i(\sigma_{-i},\sigma_{i})\geq g_i(\sigma_{-i},\sigma'_{i})
$$

\section{Etude de l'équilibre de Nash en stratégies mixtes}
Nous allons nous poser (et y répondre) les questions suivantes sur
les EN en stratégies mixtes.\\%
\begin{enumerate}
\item Quel rapport existe-t-il entre les EN en stratégies et en stratégies pures ? %
\item Comment calculer les EN en stratégies mixtes ? %
\item Quel est le rapport entre les EN en stratégies mixtes et l'élimination itérée de stratégies dominées ?%
\item Existence des EN en stratégies mixtes ?
\end{enumerate}%

\subsection{Équilibres de Nash en stratégies pures et mixtes}

\begin{proposition} $\sigma$ est un \'equilibre de Nash en
strat\'egies mixtes si et seulement si :
$$
\forall i,~~\forall s_i\in S_i~~~~~~~
g_i(\sigma_{-i},\sigma_{i})\geq g_i(\sigma_{-i},s_{i})
$$
\end{proposition} La partie ``seulement si'' provient du fait
qu'une strat\'egie pure est aussi une strat\'egie mixte, la partie
``si'' de la lin\'earit\'e de $g_i$ en $\sigma_i$.

\begin{exemple} Chercher un \'equilibre de Nash du jeu de Pile ou Face, puis
du jeu de ``Pierre, Papier, Ciseaux". %
\end{exemple}

\begin{proposition} Tout EN en strat\'egies pures est aussi un EN
en strat\'egies mixtes.
\end{proposition}
{\bf Preuve }: Corollaire de la proposition pr\'ec\'edente.

Par conséquent, considérer les stratégies mixtes nous donne plus
d'équilibres qu'avec les stratégies pures. Pour les équilibres ne
faisant intervenir que des stratégies pures, on parle d'équilibres
de Nash en stratégies pures.

\subsection{Recherche des équilibres en stratégies mixtes}
\begin{definition}
Dans un jeu fini, le support de $\sigma_i\in \Sigma_i$ est
$$
supp(\sigma_i)=\{s_i\in S_i, \sigma_i(s_i)>0\}
$$
\end{definition}

\begin{proposition}(indifférence sur le support)
Si $\sigma$ est un \'equilibre de Nash en strat\'egies mixtes,
$$
\forall i, ~~\forall s_i,s'_i\in
supp(\sigma_i)~~~~~g_i(\sigma_{-i},s_{i})= g_i(\sigma_{-i},s'_{i})
$$
\end{proposition}
En effet, si une stratégie du support donnait un meilleur paiement
que $\sigma_i$, elle constituerait une déviation profitable. Elles
donnent donc toutes un paiement inférieur ou égal à celui de
$\sigma_i$. Mais alors, si une d'entre elles donnait un paiement
strictement inférieur à celui de $\sigma_i$, on obtient une
contradiction (utiliser le fait que cette stratégie est jouée avec
probabilité $> 0$ sous $\sigma_i$ et la linéarité de $g_i$ en
$\sigma_i$).

\begin{proposition} $\sigma$ est un \'equilibre de Nash en
strat\'egies mixtes si et seulement si pour tout $i$, %
\begin{itemize}%
\item $\forall s_i\in supp(\sigma_i)$
$$
g_i(\sigma_{-i},s_{i})= g_i(\sigma_{-i},\sigma_{i})
$$
\item $\forall s_i\not\in supp(\sigma_i)$
$$
g_i(\sigma_{-i},s_{i})\leq g_i(\sigma_{-i},\sigma_{i})
$$
\end{itemize}%
\end{proposition}
Conséquence des résultats précédents.

Ceci fournit une proc\'edure de recherche des \'equilibres :
\begin{enumerate}
\item essayer tous les supports possibles ;%
\item résoudre les probabilités rendant chaque joueur
indifférent sur son
support ;%
\item vérifier que les stratégies hors du support ne donnent pas
un paiement supérieur.
\end{enumerate}

{\bf Exemple} : Quels sont les équilibres de Nash  en stratégies
pures et mixtes du jeu suivant ?
\begin{center}
\begin{game}{2}{2}
     &  $G$    &  $D$ \\
$H$ & $1,1$    & $0,4$\\
$B$ & $0,2$    & $2,1$\\
\end{game}\hspace*{20mm}%
\end{center}
Même question avec le jeu de "Pierre, Papier, Ciseaux".

\subsection{Stratégies mixtes et dominées}

Nous commençons par étendre la relation de domination stricte aux
stratégies mixtes.
\begin{definition}
On dit que  $\sigma_i$ domine $\sigma'_i$ (strictement) lorsque
pour toute $\sigma_{-i}$,
$$
g_i(\sigma_i,\sigma_{-i}) > g_i(\sigma'_i,\sigma_{-i})
$$
\end{definition}

Pour vérifier qu'une stratégie en domine une autre, il suffit de
le vérifier face à tout profil de stratégies pures.
\begin{proposition}
$\sigma_i$ domine $\sigma'_i$ si et seulement si pour tout
$s_{-i}$,
$$
g_i(\sigma_i,s_{-i}) > g_i(\sigma'_i,s_{-i})
$$
\end{proposition}
La condition nécessaire est évidente. Pour la condition
suffisante, utiliser la linéarité.


Une strat\'egie non domin\'ee par une strat\'egie pure peut \^etre
domin\'ee par une strat\'egie mixte.

\begin{exemple}
Dans le jeu suivant :
\begin{center}
\begin{game}{3}{3}
     &  $G$  &  $M$  & $D$\\
$h$ & $1,1$  & $0,2$ & $0,4$\\
$m$ & $0,2$  & $5,0$ & $1,6$\\
$b$ & $0,2$  & $1,1$ & $2,1$\\
\end{game}\hspace*{20mm}%
\end{center}
$M$ est strictement domin\'e par $\frac12G+\frac12D$.%
\end{exemple}

Nous pouvons redéfinir la procédure d'élimination itérée de
stratégies dominées en prenant en compte les stratégies dominées
par des stratégies mixtes.
\begin{itemize}
\item {\bf \'Etape 1 :} $S_i^0 = S_i$.%
\item {\bf \'Etape 2 :} %
$$
S_i^1 = \{\mbox{elts.\,de $S_i$ non dominés par des elts.\,de
$\Sigma_i$}\}
$$%
$$
\Sigma_i^1 = \{\mbox{stratégies mixtes à support dans $S_i^1$}\}
$$
\item {\bf \'Etape $k+1$ :}%
$$%
S_i^{k+1} = \{s_i \in S_i^k, \not\exists\sigma_i \in \Sigma^{k}_i,
\forall s_{-i}\in
S^{k}_{-i}, g_i(\sigma_i,s_{-i})>g_i(s_i,s_{-i})\}$$%
Stratégies non dominées {\em face aux stratégies de } $S^k_i$%
$$
\Sigma_i^k = \{\mbox{stratégies mixtes à support dans $S_i^k$}\}
$$
\item {\bf \'Etape $\infty$ :} $S_i^{\infty} = \cap_k S^k_i$.
\end{itemize}

Les équilibres de Nash en stratégies mixtes résistent à la
procédure d'élimination itérée de stratégies (strictement)
dominées.

\begin{proposition}
Si $\sigma$ est un EN en stratégies mixtes:
$$
\forall i\ \  supp(\sigma_i) \in S_i^\infty
$$
\end{proposition}
Pour la démonstration, suivre le même raisonnement que pour
l'élimination des stratégies dominées par des stratégies pures.
Plus généralement:
\begin{proposition}
Les EN en stratégies mixtes d'un jeu sont les mêmes que celui du
jeu duquel on a éliminé les stratégies strictement
dominées.
\end{proposition}

{\bf Exemple} : r\'esoudre le jeu pr\'ec\'edent.

\subsection{Existence des équilibres de Nash en stratégies mixtes}
On a un équilibre de Nash lorsque chaque joueur joue optimalement
face aux stratégies des autres joueurs.  Pour caractériser ces
équilibres, nous sommes intéressés par l'application qui associe à
chaque $\sigma_{-i}$ l'ensemble des stratégies $\sigma_i$ qui
donnent le meilleur paiement face à $\sigma_{-i}$. Cette
correspondance s'appelle la correspondance de meilleures réponses.

\begin{definition}
Une correspondance de $A$ vers $B$ est une application de $A$ vers
les parties de $B$.%
\end{definition}

\begin{definition}
 La correspondance de meilleure r\'eponse du joueur $i$ est la correspondance de
$\Sigma_{-i}$ vers $\Sigma_i$ donnée par :
$$
MR_i(\sigma_{-i})=\arg\max_{\sigma_i\in
\Sigma_i}g_i(\sigma_{-i},\sigma_i)
$$
La correspondance de meilleures réponses du jeu $G$ est la
correspondance de $\Sigma$ vers $\Sigma$ donnée par
$$
MR(\sigma) = \Pi_i MR_i(\sigma_{-i})
$$
\end{definition}
Les meilleures réponses permettent de caractériser les équilibres
de Nash :

\begin{proposition}
$\sigma$ est un EN si et seulement si :
$$
\forall i\ \  \sigma_i \in MR_i(\sigma_{-i})
$$
i.e.
$$
\sigma \in MR(\sigma)
$$
\end{proposition}

On a les propriétés suivantes de la correspondance de meilleures
réponses :

\begin{proposition}
\begin{itemize}
\item Pour tout $\sigma$, $MR(\sigma)$ est non vide%
\item Pour tout $\sigma$, $MR(\sigma)$ est convexe%
\item Si tous les $S_i$ sont finis, le graphe de $MR$ est fermé,
i.e.\,si $\sigma_n\to \sigma$, $\tau_n\to \tau$ et $\tau_n\in
MR(\sigma_n)$ alors $\tau\in MR(\sigma)$.%
\end{itemize}
\end{proposition}

On rappelle le théorème d'existence de point fixe d'une
correspondance de Kakutani :
\begin{theoreme}[de Kakutani] Soit $X$
un compact convexe non vide de $\reals^n$, et $F$ une
correspondance de $X$ vers $X$ telle que :
\begin{itemize}
\item $\forall x\in X$, $F(x)$ est convexe et non vide ; \item Le
graphe de $F$ est ferm\'e.
\end{itemize}
Alors $F$ admet un point fixe, i.e.\, il existe $x\in X$ tel que
$x\in F(x)$. %
\end{theoreme}

{\bf Exercice} Montrer que les hypothèses
\begin{itemize}
\item $X$ compact ;%
\item $X$ convexe ; \item $F(x)$ non vide pour
tout $x$ ;%
\item $F(x)$ convexe pour tout $x$ ;%
\item le graphe de $F$ est ferm\'e
\end{itemize}
sont nécessaires.

En corollaire du théorème précédent, on a le théorème d'existence
d'équilibres de Nash en stratégies mixtes :


\begin{theoreme}[Glicksberg, Nash] Tout jeu fini admet un
équilibre de Nash en stratégies mixtes.
\end{theoreme}

\section{Le jeu de p\'enalties}

Un joueur d\'ecide de tirer soit \`a droite, soit \`a gauche du
gardien. Le gardien peut plonger soit \`a droite, soit \`a gauche.
Le tir peut \^etre bien tir\'e ou sortir complètement (avec
probabilit\'e $p_D$ pour un tir \`a droite, $p_G$ \`a gauche). Le
gardien arr\^ete toujours la balle s'il plonge du bon c\^ot\'e. Il
y a donc but si :
\begin{enumerate}
\item La balle est bien tir\'ee ; %
\item Le gardien plonge du mauvais c\^ot\'e.
\end{enumerate}

On demande de résoudre les questions suivantes : L'objectif du
joueur est de maximiser les possibilit\'es de marquer but, et le
gardien veut minimiser ces possibilit\'es.
\begin{itemize}
\item Repr\'esenter le jeu sous forme normale. %
\item Quels sont les \'equilibres de Nash ? %
\item Un droitier tire-t-il plus souvent \`a gauche ou \`a droite
? %
\item Face \`a un droitier, le gardien plonge-t-il plus souvent
\`a gauche ou \`a droite ?
\end{itemize}
Indice : Pour un droitier, $p_G > p_D$.

\section{Course au brevet} Deux entreprises se lancent dans une course au
brevet. Chacune investit un montant positif dans la recherche et
celui qui investit le plus reçoit une subvention $V>0$. En cas
d'\'egalit\'e, chacun reçoit $V/2$. On prend l'intervalle $\lbrack
0,V\rbrack$ comme ensemble de strat\'egies pour chaque joueur.

\begin{enumerate}
\item Montrer que ce jeu n'a pas d'\'equilibre (en strat\'egies
pures). %
\item Montrer qu'il existe un \'equilibre en strat\'egies mixtes
dans lequel joue une loi \`a support ``plein'' $\lbrack
0,V\rbrack$.\
\end{enumerate}
